% =====================================================================
\chapter{Códigos de fuente y código de Huffman}
\label{tema:3}
% =====================================================================

Este tema da a la entropía su significado operativo: es la longitud media mínima, en bits por símbolo, de una descripción sin pérdidas de la fuente. Se estudian los códigos prefijo, la desigualdad de Kraft, las cotas para la longitud óptima y el algoritmo de Huffman~\cite{huffman}. Se sigue~\cite[cap.~5]{ct}, en particular §5.1--5.4, §5.6 y §5.8; MacKay~\cite[cap.~5]{mackay} presenta el mismo material con el nombre de \emph{códigos de símbolos}, que es también el enfoque del tema «Source coding and data compression» de~\cite{camit}.

\section{Códigos de símbolos}

\begin{definicion}[Código]\label{def:codigo}
Un \emph{código binario} para una variable $X$ con alfabeto $\X$ es una aplicación $C\colon\X\to\{\bin{0},\bin{1}\}^{*}$ que asigna a cada símbolo una cadena finita de bits, su \emph{palabra de código}, de longitud $\ell(x)$. La \emph{extensión} de $C$ codifica una secuencia concatenando palabras: $C(x_1\cdots x_n)=C(x_1)\cdots C(x_n)$. La \emph{longitud media} del código es $L(C)=\sum_{x}p(x)\,\ell(x)$.
\end{definicion}

\begin{definicion}[Clases de códigos]\label{def:clases}
Un código es
\begin{enumerate}
  \item \emph{no singular} si $C$ es inyectiva;
  \item \emph{unívocamente descodificable} si su extensión es inyectiva, es decir, si toda cadena de bits procede de, a lo sumo, una secuencia de símbolos;
  \item \emph{prefijo} o \emph{instantáneo} si ninguna palabra de código es prefijo de otra distinta.
\end{enumerate}
\end{definicion}

Cada clase contiene a la siguiente: todo código prefijo es unívocamente descodificable, porque el receptor reconoce el final de cada palabra en cuanto la lee completa, sin mirar los bits posteriores; y todo código unívocamente descodificable es no singular.

\addtocounter{definicion}{1}%

\section{La desigualdad de Kraft}

A cada cadena binaria $c=c_1\cdots c_\ell$ se le asocia el \emph{intervalo diádico} $J(c)=[0.c_1\cdots c_\ell,\ 0.c_1\cdots c_\ell+2^{-\ell})\subseteq[0,1)$, donde $0.c_1\cdots c_\ell=\sum_{i}c_i2^{-i}$. Es el conjunto de números de $[0,1)$ cuya expansión binaria empieza por $c$, y tiene longitud $2^{-\ell}$. Dos intervalos diádicos están contenidos uno en otro o son disjuntos, y $J(c')\subseteq J(c)$ si y solo si $c$ es prefijo de $c'$.

\begin{teorema}[Desigualdad de Kraft]\label{teo:kraft}
Las longitudes $\ell_1,\dots,\ell_m$ de las palabras de un código binario prefijo cumplen
\[
\sum_{i=1}^{m}2^{-\ell_i}\le 1.
\]
Recíprocamente, si unos enteros positivos $\ell_1,\dots,\ell_m$ cumplen esta desigualdad, existe un código prefijo con esas longitudes.
\end{teorema}

\addtocounter{figure}{1}%

\begin{observacion}
McMillan demostró que la desigualdad de Kraft la cumplen también todos los códigos unívocamente descodificables~\cite[§5.5]{ct}. Por tanto, restringirse a códigos prefijo no tiene coste: cualquier conjunto de longitudes alcanzable con un código unívocamente descodificable lo es también con uno prefijo.
\end{observacion}

\section{Cotas para la longitud óptima}

\begin{teorema}[Cota inferior]\label{teo:cota-inferior}
La longitud media de todo código prefijo para $X$ cumple $L\ge H(X)$, con igualdad si y solo si $p(x)=2^{-\ell(x)}$ para todo $x$.
\end{teorema}

La demostración interpreta cada código prefijo como un modelo probabilístico implícito, $r(x)\propto2^{-\ell(x)}$: un código asigna palabras cortas a lo que «cree» probable. Esta lectura es la mitad de la equivalencia entre predecir y comprimir que se completa en el \cref{tema:4}.

\begin{definicion}[Código de Shannon]\label{def:shannon}
Un \emph{código de Shannon} para $p$ es un código prefijo con longitudes $\ell(x)=\lceil-\log p(x)\rceil$.
\end{definicion}

Estas longitudes cumplen la desigualdad de Kraft, porque $2^{-\lceil-\log p(x)\rceil}\le2^{\log p(x)}=p(x)$ y las $p(x)$ suman~1; por el \cref{teo:kraft}, el código existe.

\begin{teorema}[Longitud óptima]\label{teo:cotas}
Sea $L^{*}$ la mínima longitud media de un código prefijo para $X$. Entonces
\[
H(X)\le L^{*}<H(X)+1.
\]
\end{teorema}

El exceso de hasta un bit se debe a que las longitudes han de ser enteras. Se reparte entre muchos símbolos codificando bloques.

\begin{corolario}[Codificación por bloques]\label{cor:bloques}
Sean $X_1,\dots,X_n$ independientes y con la misma distribución que $X$, y sea $L_n^{*}$ la longitud media mínima de un código prefijo para el bloque $(X_1,\dots,X_n)$. Entonces
\[
H(X)\le\frac{L_n^{*}}{n}<H(X)+\frac1n.
\]
En particular, $L_n^{*}/n\to H(X)$ cuando $n\to\infty$.
\end{corolario}

El \cref{cor:bloques} es una forma del \emph{teorema de codificación de fuente} de Shannon: la entropía es el número mínimo de bits por símbolo con que puede describirse la fuente sin pérdidas, y ese mínimo se alcanza asintóticamente~\cite[cap.~4]{mackay}.

\section{El código de Huffman}

\begin{definicion}[Algoritmo de Huffman]\label{def:huffman}
Dada una distribución sobre $K\ge 2$ símbolos, el \emph{algoritmo de Huffman} construye un árbol binario: mientras quede más de un nodo, fusiona los dos nodos de menor probabilidad en un nodo nuevo cuya probabilidad es la suma de ambas. Al terminar, las ramas que salen de cada nodo se etiquetan con \bin{0} y \bin{1}, y la palabra de cada símbolo es la secuencia de etiquetas desde la raíz hasta su hoja.
\end{definicion}

Cada fusión añade un bit a la palabra de todos los símbolos de los dos nodos fusionados; por eso los símbolos poco probables, que se fusionan pronto, reciben palabras largas.

\addtocounter{definicion}{1}%

\addtocounter{figure}{1}%

La optimalidad del algoritmo se apoya en dos propiedades de los códigos óptimos.

\begin{lema}\label{lem:optimo}
Sea $p_1\ge p_2\ge\dots\ge p_K>0$, con $K\ge 2$. Existe un código prefijo óptimo con longitudes $\ell_1,\dots,\ell_K$ tal que
\begin{enumerate}
  \item $\ell_1\le\ell_2\le\dots\le\ell_K$;
  \item las palabras de los dos símbolos menos probables, $K-1$ y $K$, tienen la longitud máxima $\ell_K$ y solo difieren en el último bit.
\end{enumerate}
\end{lema}

\begin{teorema}[Optimalidad del código de Huffman]\label{teo:huffman}
El código de Huffman tiene longitud media mínima entre todos los códigos prefijo para $p$.
\end{teorema}

Por los Teoremas~\ref{teo:cotas} y~\ref{teo:huffman}, $H(p)\le L_{\text{Huffman}}<H(p)+1$, con $L_{\text{Huffman}}=H(p)$ cuando todas las probabilidades son potencias de $1/2$. El código óptimo no es único: con otros desempates en el \cref{ej:huffman} se obtiene un árbol distinto (\cref{eje:3-dos-huffman}).

\section{El límite de los códigos de símbolos}

Un código de símbolos gasta un número entero de bits por símbolo, y al menos un bit. Si $p=(0{,}9;\ 0{,}1)$, la entropía es $H_b(0{,}1)=0{,}469$ bits, pero cualquier código de símbolos para dos resultados gasta 1 bit por símbolo: más del doble. La codificación por bloques del \cref{cor:bloques} reduce el exceso, pero un bloque de $n$ símbolos sobre un alfabeto de $K$ símbolos tiene $K^{n}$ valores posibles, y la tabla del código crece exponencialmente.

\addtocounter{definicion}{1}%

